\clearpage
\thispagestyle{empty}
\pdfbookmark[0]{Dioptrical Investigations. Edition and sources}{dioptrische-source-notice}
\begin{center}\Large\textbf{Dioptrical Investigations}\end{center}

AI-translated from the German text in Carl Friedrich Gauss, \textit{Werke}, volume V (Göttingen, 1867), native leaves 249--282. This edition contains the formal title and imprint, repeated heading, introduction and all 22 articles, with AI-produced typesetting. The work was presented on 10 December 1840. Its title gives the first volume of the Göttingen society's transactions and the Göttingen imprint of 1843. The example in Article 19 retains the original units and decimal commas.

Bernard Rosett's English typescript \textit{Optical Investigations}, dated July 1941, was consulted for its title page, translator's note, introduction and Articles 1--22. Rosett also wrote the separate appendix. An opening excerpt of the English abstract in Richard Taylor (ed.), \textit{Scientific Memoirs}, vol.~III (1843), Art.~XIX, pp.~490--498, communicated by W.~H.~Miller, was also consulted. The Karlsruhe Institute of Technology's articles \textit{Brechung} and \textit{Linsen} provided modern German optics terminology.

\noindent\href{https://gdz.sub.uni-goettingen.de/id/PPN236006339}{German source (1867)}\quad
\href{https://wp.optics.arizona.edu/jsasian/wp-content/uploads/sites/33/2016/02/Gauss.pdf}{Rosett (1941)}\quad
\href{https://archive.org/details/scientificmemoir03memo/page/490/mode/1up}{English abstract (1843)}\quad
\par\smallskip\noindent\href{https://labs.physik.kit.edu/prakt-klass-physik-manuals/P1/Geometrische_Optik/doc/Hinweise-Brechung/}{KIT: Brechung}\quad
\href{https://labs.physik.kit.edu/prakt-klass-physik-manuals/P1/Geometrische_Optik/doc/Hinweise-Linsen/}{KIT: Linsen}
\clearpage
